% =========================================================
% Section: Introduction to Lambda Calculus
% Chapter: Lambda Calculus — Volume I
% =========================================================

\begin{tcolorbox}[colback=gray!6, colframe=gray!40, arc=2pt,
  left=6pt, right=6pt, top=4pt, bottom=4pt,
  title={\small\textbf{Section Toolkit --- Quick Reference}},
  fonttitle=\small\bfseries]
\small
\begin{tabular}{l l}
\toprule
\textbf{Concept} & \textbf{One-line summary} \\
\midrule
$\lambda$-term
  & Variable $x$, abstraction $\lambda x.M$, application $MN$ \\
$\alpha$-equivalence
  & $\lambda x.M =_\alpha \lambda y.M[y/x]$ when $y$ fresh \\
$\beta$-reduction
  & $(\lambda x.M)N \to_\beta M[N/x]$ \\
Normal form
  & Term with no $\beta$-redex; may not exist (non-termination) \\
Church boolean \textbf{T}
  & $\lambda x.\lambda y.x$ \\
Church boolean \textbf{F}
  & $\lambda x.\lambda y.y$ \\
Church numeral $\overline{n}$
  & $\lambda f.\lambda x.f^n(x)$ ($f$ applied $n$ times to $x$) \\
Type $\tau \to \sigma$
  & Function type: input of type $\tau$, output of type $\sigma$ \\
Typing judgment
  & $\Gamma\vdash M : \tau$ reads ``$M$ has type $\tau$ in context $\Gamma$'' \\
Curry--Howard
  & Propositions are types; proofs are terms; $A\to B$ is implication \\
\bottomrule
\end{tabular}
\end{tcolorbox}

% ---------------------------------------------------------
% Part I: Untyped Lambda Calculus
% ---------------------------------------------------------

\subsection*{Untyped Lambda Calculus}

\begin{tcolorbox}[colback=propbox, colframe=propborder, arc=2pt,
  left=6pt, right=6pt, top=4pt, bottom=4pt,
  title={\small\textbf{Definition (Lambda Terms)}},
  fonttitle=\small\bfseries]
\label{def:lambda-terms}
Fix a countably infinite set $\mathcal{V}$ of \textbf{variables}.
The set $\Lambda$ of \textbf{lambda terms} is defined inductively:
\begin{enumerate}
  \item[(i)] (\textit{Variable}) Every $x \in \mathcal{V}$ is a
             lambda term.
  \item[(ii)] (\textit{Abstraction}) If $x \in \mathcal{V}$ and
              $M \in \Lambda$, then $(\lambda x.M) \in \Lambda$.
  \item[(iii)] (\textit{Application}) If $M, N \in \Lambda$,
               then $(MN) \in \Lambda$.
\end{enumerate}
Application is left-associative: $MNP := (MN)P$.
Abstraction extends as far right as possible:
$\lambda x.MN := \lambda x.(MN)$.
\end{tcolorbox}

\begin{remark*}[Fully quantified form]
$\Lambda$ is the smallest set satisfying (i)--(iii).
\end{remark*}

\begin{remark*}[Flash]
\FlashQ{What are the three syntactic forms of a lambda term?}
\FlashA{Variable $x$; abstraction $\lambda x.M$ (a function with
  parameter $x$ and body $M$); application $MN$ (apply $M$ to $N$).}
\FlashHint{The entire expressiveness of lambda calculus --- booleans,
  numbers, recursion, data structures --- is encoded using only these
  three forms.}
\end{remark*}

\begin{tcolorbox}[colback=propbox, colframe=propborder, arc=2pt,
  left=6pt, right=6pt, top=4pt, bottom=4pt,
  title={\small\textbf{Definition (Free and Bound Variables)}},
  fonttitle=\small\bfseries]
\label{def:free-bound-variables}
The sets $\mathrm{FV}(M)$ and $\mathrm{BV}(M)$ of
\textbf{free} and \textbf{bound} variables of $M \in \Lambda$
are defined by structural recursion:
\begin{align*}
\mathrm{FV}(x)        &= \{x\},
  & \mathrm{BV}(x)        &= \emptyset \\
\mathrm{FV}(\lambda x.M) &= \mathrm{FV}(M) \setminus \{x\},
  & \mathrm{BV}(\lambda x.M) &= \mathrm{BV}(M) \cup \{x\} \\
\mathrm{FV}(MN)       &= \mathrm{FV}(M) \cup \mathrm{FV}(N),
  & \mathrm{BV}(MN)       &= \mathrm{BV}(M) \cup \mathrm{BV}(N)
\end{align*}
A term $M$ is \textbf{closed} (a \textbf{combinator}) if
$\mathrm{FV}(M) = \emptyset$.
\end{tcolorbox}

\begin{remark*}[Fully quantified form]
$\mathrm{FV}$ and $\mathrm{BV}$ are total functions $\Lambda\to\mathcal{P}(\mathcal{V})$
defined by the equations above.
\end{remark*}

\begin{remark*}[Flash]
\FlashQ{What makes a variable free vs.\ bound in a lambda term?}
\FlashA{A variable occurrence is \emph{bound} if it is within the scope
  of a $\lambda$ that binds it; \emph{free} otherwise.
  In $\lambda x.(xy)$: $x$ is bound, $y$ is free.
  A closed term has no free variables.}
\FlashHint{The binder $\lambda x$ captures every free occurrence of $x$
  in its body.  Shadowing: inner $\lambda x$ captures occurrences in
  its body even if outer $\lambda x$ exists.}
\end{remark*}

\begin{tcolorbox}[colback=propbox, colframe=propborder, arc=2pt,
  left=6pt, right=6pt, top=4pt, bottom=4pt,
  title={\small\textbf{Definition ($\alpha$-Equivalence)}},
  fonttitle=\small\bfseries]
\label{def:alpha-equivalence}
Two lambda terms are \textbf{$\alpha$-equivalent}, written
$M =_\alpha N$, if one can be obtained from the other by
consistently renaming bound variables.  Formally,
$=_\alpha$ is the smallest congruence on $\Lambda$ satisfying:
\[
  \lambda x.M =_\alpha \lambda y.M[y/x]
  \quad \text{whenever } y \notin \mathrm{FV}(M).
\]
We work with $\alpha$-equivalence classes throughout: terms
that differ only in bound variable names are identified.
\end{tcolorbox}

\begin{remark*}[Fully quantified form]
$=_\alpha$ is an equivalence relation on $\Lambda$, and
$\lambda x.M =_\alpha \lambda y.M[y/x]$ for all $y\notin\mathrm{FV}(M)$.
\end{remark*}

\begin{remark*}[Flash]
\FlashQ{What is $\alpha$-equivalence and why do we need it?}
\FlashA{$\alpha$-equivalence identifies terms differing only in bound
  variable names: $\lambda x.x =_\alpha \lambda y.y$.
  We need it because the \emph{name} of a bound variable carries no
  information --- only the binding structure matters.
  It prevents variable capture in substitution.}
\FlashHint{Think of it as ``the name of a local variable in a function
  doesn't matter.'' $f(x) = x^2$ and $f(t) = t^2$ define the same function.}
\end{remark*}

\begin{tcolorbox}[colback=propbox, colframe=propborder, arc=2pt,
  left=6pt, right=6pt, top=4pt, bottom=4pt,
  title={\small\textbf{Definition ($\beta$-Reduction)}},
  fonttitle=\small\bfseries]
\label{def:beta-reduction}
The \textbf{$\beta$-reduction} relation $\to_\beta$ on $\Lambda$
is the smallest compatible relation satisfying:
\[
  (\lambda x.M)N \to_\beta M[N/x]
\]
where $M[N/x]$ denotes the capture-avoiding substitution of $N$
for all free occurrences of $x$ in $M$.
The reflexive-transitive closure is written $\twoheadrightarrow_\beta$.
A term with no $\beta$-redex (no subterm of the form
$(\lambda x.M)N$) is in \textbf{$\beta$-normal form}.
\end{tcolorbox}

\begin{remark*}[Fully quantified form]
$\to_\beta$ is a relation on $\Lambda/\!=_\alpha$ satisfying the
one-step reduction axiom above, closed under abstraction and
application contexts.
\end{remark*}

\begin{remark*}[Flash]
\FlashQ{State the $\beta$-reduction rule and give an example.}
\FlashA{$(\lambda x.M)N \to_\beta M[N/x]$: apply the function
  $\lambda x.M$ to argument $N$ by substituting $N$ for $x$ in $M$.
  Example: $(\lambda x.x\,x)(\lambda y.y)
  \to_\beta (\lambda y.y)(\lambda y.y)
  \to_\beta \lambda y.y$.}
\FlashHint{``Capture-avoiding'' means if $N$ contains a free variable
  $z$ and $M$ contains $\lambda z.(...)$, rename the bound $z$ before
  substituting.  $\alpha$-equivalence handles this.}
\end{remark*}

% ---------------------------------------------------------
% Church Encodings
% ---------------------------------------------------------

\subsection*{Church Encodings}

\begin{tcolorbox}[colback=propbox, colframe=propborder, arc=2pt,
  left=6pt, right=6pt, top=4pt, bottom=4pt,
  title={\small\textbf{Definition (Church Booleans)}},
  fonttitle=\small\bfseries]
\label{def:church-booleans}
The \textbf{Church booleans} are the closed lambda terms:
\[
  \mathbf{T} := \lambda x.\lambda y.\,x
  \qquad
  \mathbf{F} := \lambda x.\lambda y.\,y
\]
The \textbf{if-then-else} combinator is:
\[
  \mathbf{if} := \lambda b.\lambda t.\lambda f.\,b\,t\,f
\]
\end{tcolorbox}

\begin{remark*}[Fully quantified form]
$\mathbf{T}, \mathbf{F}, \mathbf{if} \in \Lambda$ are closed terms.
$\mathbf{if}\,\mathbf{T}\,M\,N \twoheadrightarrow_\beta M$ and
$\mathbf{if}\,\mathbf{F}\,M\,N \twoheadrightarrow_\beta N$
for all $M,N\in\Lambda$.
\end{remark*}

\begin{remark*}[Flash]
\FlashQ{What are the Church booleans and how does if-then-else work?}
\FlashA{$\mathbf{T} = \lambda x.\lambda y.x$ (select first);
  $\mathbf{F} = \lambda x.\lambda y.y$ (select second).
  $\mathbf{if}\,\mathbf{T}\,M\,N
  = (\lambda b\lambda t\lambda f.btf)\,\mathbf{T}\,M\,N
  \twoheadrightarrow_\beta \mathbf{T}\,M\,N
  = (\lambda x\lambda y.x)\,M\,N
  \twoheadrightarrow_\beta M$.}
\FlashHint{The key insight: a boolean \emph{is} a binary selector.
  True selects the first branch; false selects the second.
  No primitive boolean type is needed.}
\end{remark*}

\begin{remark*}[Intuition]
Every Church encoding works by the same principle: data is represented
as a function that \emph{performs} the relevant case analysis.  A boolean
selects one of two branches.  A natural number applies a function $n$
times.  A list folds over its elements.  This is not a trick --- it is
the deepest fact about computation: data and control flow are the same thing.
\end{remark*}

\begin{tcolorbox}[colback=propbox, colframe=propborder, arc=2pt,
  left=6pt, right=6pt, top=4pt, bottom=4pt,
  title={\small\textbf{Definition (Church Numerals)}},
  fonttitle=\small\bfseries]
\label{def:church-numerals}
The \textbf{Church numeral} $\overline{n}$ for $n \in \mathbb{N}$
is the closed lambda term:
\[
  \overline{n} := \lambda f.\lambda x.\,f^n(x)
\]
where $f^0(x) := x$ and $f^{n+1}(x) := f(f^n(x))$.
Explicitly: $\overline{0} = \lambda f.\lambda x.\,x$,\quad
$\overline{1} = \lambda f.\lambda x.\,f\,x$,\quad
$\overline{2} = \lambda f.\lambda x.\,f(f\,x)$.
\end{tcolorbox}

\begin{remark*}[Fully quantified form]
$\forall n\in\mathbb{N}$: $\overline{n}\in\Lambda$ is closed and
$\overline{n}\,f\,x \twoheadrightarrow_\beta f^n(x)$ for all $f,x$.
\end{remark*}

\begin{remark*}[Flash]
\FlashQ{What is the Church numeral $\overline{n}$?}
\FlashA{$\overline{n} = \lambda f.\lambda x.f^n(x)$: a function that
  takes $f$ and $x$ and applies $f$ to $x$ exactly $n$ times.
  $\overline{0}\,f\,x = x$; $\overline{1}\,f\,x = f\,x$;
  $\overline{2}\,f\,x = f(f\,x)$.}
\FlashHint{The successor function is
  $\mathbf{S} = \lambda n.\lambda f.\lambda x.f\,(n\,f\,x)$:
  apply $f$ once more than $n$ does.}
\end{remark*}

% ---------------------------------------------------------
% Part II: Simply Typed Lambda Calculus
% ---------------------------------------------------------

\subsection*{Simply Typed Lambda Calculus ($\lambda^\to$)}

\begin{tcolorbox}[colback=propbox, colframe=propborder, arc=2pt,
  left=6pt, right=6pt, top=4pt, bottom=4pt,
  title={\small\textbf{Definition (Simple Types)}},
  fonttitle=\small\bfseries]
\label{def:simple-types}
Fix a set of \textbf{base types} (e.g.\ $o$ for propositions).
The set $\mathcal{T}$ of \textbf{simple types} is:
\begin{enumerate}
  \item[(i)] Every base type is a simple type.
  \item[(ii)] If $\sigma, \tau \in \mathcal{T}$ then
              $(\sigma \to \tau) \in \mathcal{T}$.
\end{enumerate}
$\to$ is right-associative: $\sigma\to\tau\to\rho := \sigma\to(\tau\to\rho)$.
\end{tcolorbox}

\begin{remark*}[Fully quantified form]
$\mathcal{T}$ is the smallest set of expressions satisfying (i) and (ii).
\end{remark*}

\begin{remark*}[Flash]
\FlashQ{What is a simple type and how is the function type formed?}
\FlashA{Base types plus function types: if $\sigma$ and $\tau$ are types
  then $\sigma\to\tau$ is the type of functions from $\sigma$ to $\tau$.
  Right-associative: $A\to B\to C$ means $A\to(B\to C)$, a curried
  two-argument function.}
\end{remark*}

\begin{tcolorbox}[colback=propbox, colframe=propborder, arc=2pt,
  left=6pt, right=6pt, top=4pt, bottom=4pt,
  title={\small\textbf{Definition (Typing Judgment)}},
  fonttitle=\small\bfseries]
\label{def:typing-judgment}
A \textbf{context} $\Gamma$ is a finite list of pairs
$x_1:\sigma_1,\ldots,x_n:\sigma_n$ with distinct variable names.
The \textbf{typing judgment} $\Gamma \vdash M : \tau$ is defined
by the following rules:
\[
  \frac{x:\tau \in \Gamma}{\Gamma \vdash x : \tau}
  \qquad
  \frac{\Gamma, x:\sigma \vdash M : \tau}
       {\Gamma \vdash \lambda x.M : \sigma\to\tau}
  \qquad
  \frac{\Gamma \vdash M : \sigma\to\tau \quad \Gamma \vdash N : \sigma}
       {\Gamma \vdash MN : \tau}
\]
\end{tcolorbox}

\begin{remark*}[Fully quantified form]
$\Gamma\vdash M:\tau$ is derivable iff there exists a derivation tree
built from the three rules above.
\end{remark*}

\begin{remark*}[Flash]
\FlashQ{State the three typing rules for simply typed lambda calculus.}
\FlashA{(Var): a variable has its declared type.
  (Abs): $\lambda x.M$ has type $\sigma\to\tau$ if $M$ has type $\tau$
  in a context extended with $x:\sigma$.
  (App): applying a function of type $\sigma\to\tau$ to an argument
  of type $\sigma$ gives type $\tau$.}
\FlashHint{These three rules are the computational counterpart of three
  logical rules: hypothesis, implication introduction, implication
  elimination (modus ponens).}
\end{remark*}

% ---------------------------------------------------------
% Curry-Howard Correspondence
% ---------------------------------------------------------

\subsection*{The Curry--Howard Correspondence}

\begin{tcolorbox}[colback=thmbox, colframe=thmborder, arc=2pt,
  left=6pt, right=6pt, top=4pt, bottom=4pt,
  title={\small\textbf{Theorem (Curry--Howard Correspondence)}},
  fonttitle=\small\bfseries]
\label{thm:curry-howard}
There is a correspondence between simply typed lambda calculus
and propositional intuitionistic logic:
\[
  \begin{array}{ccc}
  \textbf{Type theory} & & \textbf{Logic} \\
  \hline
  \text{type } \tau & \longleftrightarrow & \text{proposition } \tau \\
  \text{term } M : \tau & \longleftrightarrow & \text{proof of } \tau \\
  \text{function type } \sigma\to\tau & \longleftrightarrow
    & \text{implication } \sigma\Rightarrow\tau \\
  \text{abstraction } \lambda x.M & \longleftrightarrow
    & \text{implication introduction} \\
  \text{application } MN & \longleftrightarrow
    & \text{modus ponens} \\
  \text{type checking} & \longleftrightarrow & \text{proof checking}
  \end{array}
\]
\end{tcolorbox}

\begin{remark*}[Fully quantified form]
$\Gamma\vdash M:\tau$ is derivable in $\lambda^\to$ iff
$\Gamma$ (viewed as a list of hypotheses) proves $\tau$
(viewed as a proposition) in intuitionistic propositional logic.
\end{remark*}

\begin{remark*}[Flash]
\FlashQ{State the Curry--Howard correspondence in one sentence.}
\FlashA{Propositions are types, proofs are programs (lambda terms),
  and checking a proof is type-checking a term.}
\FlashHint{The simplest instance: the identity function $\lambda x.x$
  has type $A\to A$ and is simultaneously the proof of $A\Rightarrow A$.
  In Lean: \texttt{fun x => x} proves any goal of the form \texttt{A \to\ A}.}
\end{remark*}

\begin{remark*}[Intuition]
This is not a metaphor.  Every Lean tactic proof is constructing a lambda
term behind the scenes.  When you write \texttt{intro h} you are writing
$\lambda h.$; when you write \texttt{exact h} you are supplying the term
$h$.  The \texttt{\#check} command in Lean shows the type of a term ---
which, under Curry--Howard, is the proposition it proves.  Understanding
lambda calculus makes the Lean kernel transparent.
\end{remark*}

\begin{remark*}[Forward reference]
A full treatment of dependent types (where types can depend on terms),
the Calculus of Inductive Constructions, and Lean's underlying type
theory (based on the Calculus of Constructions with universes) is
deferred until sufficient formalization experience accumulates.
The present section provides the conceptual foundation; the full
machinery will be added as a later section of this chapter.
\end{remark*}
